\documentclass[11pt]{article}
% Кодировка и язык
\usepackage[utf8]{inputenc}           % Кодировка исходного файла (UTF-8)
\usepackage[T2A]{fontenc}             % Кодировка шрифтов для кириллицы
\usepackage[english,russian]{babel}   % Локализация: русские названия, переносы

% Поиск в PDF
\usepackage{cmap}                     

% Поля страницы
\usepackage{geometry}
\geometry{left=2cm, right=2cm, top=2cm, bottom=2cm}

% Математика
\usepackage{amsmath, amssymb, amsthm} % Базовые пакеты для математики
\usepackage{mathtools}                % Расширение amsmath (доп. символы и окружения)

% Картинки
\usepackage{graphicx}                 % Вставка картинок (\includegraphics)
\graphicspath{{pictures/}}            % Путь к папке с картинками
\usepackage{float}                    % Позволяет использовать [H] для жёсткого позиционирования
\usepackage{wrapfig}                  % Обтекание текста вокруг картинок

% Таблицы
\usepackage{array}                    % Расширенные возможности для таблиц
\usepackage{booktabs}                 % Красивые таблицы (\toprule, \midrule, \bottomrule)
\usepackage{multirow}                 % Объединение строк в таблицах
\usepackage{tabularx}                 % Таблицы с автоматической шириной колонок
\usepackage{ragged2e}                 % Улучшенное выравнивание (используется в tabularx)

% Списки
\usepackage{enumitem}                 % Гибкая настройка списков (itemize, enumerate)

% Колонтитулы
\usepackage{fancyhdr}                 

% Гиперссылки
\usepackage{xcolor}                   % Работа с цветами
\usepackage{hyperref}                 % Гиперссылки, оглавление, перекрёстные ссылки

% Умные ссылки
\usepackage{cleveref}                 

% Типографика
\usepackage{microtype}                % Улучшает переносы и межсловные расстояния

% Подписи к картинкам и таблицам
\usepackage{caption}                  % Настройка подписей (шрифт, отступы)

% Рисование схем и графиков
\usepackage{tikz}                     % Рисование схем, графиков, диаграмм
\usepackage{pgfplots}                 % Построение графиков функций (на основе tikz)
\pgfplotsset{compat=newest}           % Версия совместимости pgfplots

% Алгоритмы и код
\usepackage{algorithm}                % Окружения для алгоритмов
\usepackage{algorithmic}              % Синтаксис для алгоритмов
\usepackage{listings}                 % Вставка кода с подсветкой синтаксиса


% Настройка нумерации секций и подсекций
\renewcommand{\thesection}{\arabic{section}}                   
\renewcommand{\thesubsection}{\thesection.\arabic{subsection}} 


% Настройка списков: убираем лишние отступы слева
\setlist[itemize]{leftmargin=}
\setlist[enumerate]{leftmargin=}


% Определение окружений для теорем, лемм, определений и т.д.
% Нумерация привязана к номеру секции
\newtheorem{theorem}{Теорема}[section]
\newtheorem{definition}{Определение}[section]
\newtheorem{remark}{Замечание}[section]
\newtheorem{lemma}{Лемма}[section]
\newtheorem{example}{Пример}[section]
\newtheorem{consequence}{Следствие}[section]
\newtheorem{algorithmm}{Алгоритм}
\newtheorem{question}{Вопрос}[section]
\newtheorem{answer}{Ответ}[section]
\newtheorem{statement}{Утверждение}[section]


% Определение цветов
\definecolor{pblue}{rgb}{0.13,0.13,1}    
\definecolor{pgreen}{rgb}{0,0.5,0}       
\definecolor{pred}{rgb}{0.9,0,0}         
\definecolor{pgrey}{rgb}{0.46,0.45,0.48}
\newtheorem{fact}[theorem]{Факт}

\newcommand{\lectureNumber}{12}
\newcommand{\lectureTitle}{Теорема об изоморфизме. Вронскиан вектор-функций. Вронскиан скалярных функций.}
\newcommand{\courseName}{Элементы дифференциальных уравнений}
\newcommand{\semester}{3 семестр, 2026/2027 уч. год}

\pagestyle{fancy}
\fancyhf{}
\fancyhead[L]{Мишка, по телевизору тоже снег показывают!}
\fancyhead[R]{\thepage}
\setlength{\headheight}{15pt}

\hypersetup{
    colorlinks=true,
    linkcolor=blue!70!black,
    citecolor=blue!70!black,
    filecolor=blue!70!black,
    urlcolor=blue!70!black,
}

\begin{document}


\begin{titlepage}
    \centering
    
    {\large Московский физико-технический институт \par}
    {\large Высшая школа программной инженерии \par}
    \vspace{2cm}

    {\Large \textbf{\courseName} \par}
    {\large \semester \par}
    
    \vspace*{\fill}

    {\Huge \textbf{Лекция \lectureNumber} \par}
    {\Huge \textbf{\lectureTitle} \par}
    
    \vspace*{\fill}
    
    {\large Лектор: Саулин Сергей Михайлович \par}
    {\large Автор: \href{https://t.me/mallina_olga}{Фиксик ночью постарался} \par}

\end{titlepage}

\pagenumbering{arabic}

\section{А помните...}
Напомним, что мы рассматриваем однородные линейные уравнения с переменными коэффициентами
\begin{equation}
   \dot{x}= A(t)x, A(t)\in \mathrm{MAT}_{n\times n} (\mathbb{R}), t\in I \subset \mathbb{R} \text{ --- интервал}, A(t) = (a_{ij}(t))_{i,j=1}^n, a_{ij} \in C(I)  
   \label{eq:1}
\end{equation}
 

\begin{remark}
      Пусть $\mathcal{L}$ --- множество всех решений уравнения \eqref{eq:1}. Заметим, что $\mathcal{L}$ --- линейное пространство. 
\end{remark}

\begin{theorem}
    (об изоморфизме) Пусть $t_0\in I$ --- фиксированное. Тогда 
    \[ \varphi_{t_0}: \mathcal{L} \to \mathbb{R}^n, x(t)\to x(t_0)=\varphi_{t_0} x(t) \text{ является изоморфизмом.} \]
\end{theorem}

\begin{proof}
    $\forall a\in \mathbb{R}^n \ \exists! \  x(t)$ --- решение \eqref{eq:1} т.ч. $x(t_0)=a$, так как система
    $ \begin{cases}
        \dot{x} = A(t)x\\
        x(t_0)=a
    \end{cases} $
    имеет единственное решение $\Rightarrow \varphi_{t_0}$ --- биекция. \\
    Теперь докажем, что $\varphi_{t_0}$ --- гомоморфизм линейных пространств $\mathcal{L}$ и $\mathbb{R}^n$. 
    Возьмём функцию и посмотрим образ этой функции под действием $\varphi_{t_0}: \varphi_{t_0} (\lambda x(t)) = \lambda x(t_0) = \lambda \varphi_{t_0} (x(t)) $. \\ 
    Теперь, если возьмём два решения уравнения \eqref{eq:1}, то получим следующее:
    \[ \varphi_{t_0} (x_1(t)+x_2(t)) = x_1(t_0) + x_2(t_0) = \varphi_{t_0} (x_1(t)) + \varphi_{t_0}(x_2(t)) \] 
\end{proof}

\begin{consequence}
    \[ dim\, \mathcal{L} = n \] 
\end{consequence}

\begin{definition}
    ФСР уравнения \eqref{eq:1} --- произвольный базис в $\mathcal{L}$. 
\end{definition}

\begin{consequence}
    ФСР уравнения \eqref{eq:1} существует, причём $\forall$ решение уравнения \eqref{eq:1} является линейной комбинацией вектор-функций, образующих ФСР уравнения \eqref{eq:1}. 
\end{consequence}

\begin{proof}
    Пусть $e_1, \ldots e_n$ --- базис в $\mathbb{R}^n$, а $ x_1(t) = \varphi_{t_0}^{-1}(e_1), \ldots, x_n(t) = \varphi_{t_0}^{-1} (e_n) $.\\
    Тогда $x_1(t), \ldots x_n(t)$ --- базис в $\mathcal{L}$. 
\end{proof}


\section{Вронскиан вектор-функций}

Пусть $I\subset \mathbb{R}$ --- интервал, $\mathcal{F} = \{f: I\to \mathbb{R}^n , f\in C(I)$ \}.

\begin{definition}
    Пусть $f_1, \ldots, f_n \in \mathcal{F}$. Тогда вронскианом набора вектор-функций $f_1, \ldots, f_n$ называется
    \[ W_{f_1, \ldots, f_n} (t) = det\, (f_1(t) \ldots f_n(t)) \] 
    Сама матрица называется матрицей Вронского. 
\end{definition}

\begin{remark}
    Если $f_1, \ldots, f_n$ --- линейно зависимые $\Rightarrow W_{f_1, \ldots, f_n}(t) = 0 \ \forall t\in I$.\\
    В обратную сторону это утверждение неверно.
\end{remark}

\begin{example}
    Пусть $f_1(t)= 
    \begin{pmatrix}
        t^2\\
        2t  
    \end{pmatrix}
    , f_2(t) = \begin{pmatrix}
        t|t| \\
        2|t|
    \end{pmatrix}    
    $. Тогда $W_{f_1, f_2}(t) = 0$, но $f_1$ и $f_2$ линейно независимы. 
\end{example}


\begin{statement}
    Пусть $x_1(t), \ldots, x_n(t)$ --- решение уравнения \eqref{eq:1}. Тогда следующие утверждения эквиваленты:
    \begin{enumerate}
        \item $x_1, \ldots, x_n(t)$ --- линейно зависимы.
        \item $W_{x_1, \ldots, x_n}(t) \equiv 0 $ на $I$
        \item $\exists t_0\in I: W_{x_1, \ldots x_n}(t_0) = 0$
    \end{enumerate}
\end{statement}

\begin{proof}
    Нетрудно заметить (например из предыдущего замечания), что $1\Rightarrow 2 \Rightarrow 3$.\\
    Покажем $3\Rightarrow 1:$ Так как $W_{x_1,\ldots, x_n}(t_0) = 0 \Rightarrow x_1(t_0), \ldots, x_n(t_0)$ --- линейно зависимы. Тогда $\displaystyle \exists c_1, \ldots, c_n\in \mathbb{R}: \sum_{k=1}^{n} x_k(t_0)c_k = 0 $. 
    Тогда рассмотрим следующие функции: $\displaystyle x_t\equiv 0$,\\
    $ \displaystyle \tilde{x}(t)=\sum_{k=1}^{n} c_kx_k(t)$ --- решения уравнения \eqref{eq:1}, причём $x(t_0)=0$ и $\tilde{x}(t_0)=0$. 
    В силу единственности непродолжаемого решения уравнения \eqref{eq:1} $\displaystyle \sum_{k=1}^{n} c_k x_k(t) =0 \ \forall t\in I$. 
\end{proof}

\begin{theorem}
    (Лиувилля-Остроградского) Пусть $W(t)$ --- вронскиан решений $x_1(t)\ldots x_n(t)$ уравнения \eqref{eq:1}.
    Тогда  \[ \forall t_0\in I \ W(t) = W(t_0) e^{\int_{t_0}^{t} tr\,A(s)\, ds} \ \forall t\in I \] 
\end{theorem}

\begin{proof}
    Для доказательства теоремы используем следующую лемму \ref{lem:2.1}. 
    Проинтегрируем равенство из леммы:
    \[ \frac{dW}{W} = trA(t)\, dt \Rightarrow \ln \left| \frac{W(t)}{W(t_0)} \right| = \int_{t_0}^{t} trA(s)\, ds \]
\end{proof}

\begin{lemma}\label{lem:2.1}
    \[ \dot{W} = tr(A(t))W \]   
\end{lemma}

\begin{proof}
    Пусть $X(t)$ --- матрица, столбцами которой являются $x_1(t), \ldots, x_n(t)$. 
    \[ X(t) = \begin{pmatrix}
        x_1(t) \ldots x_n(t)
    \end{pmatrix} \]
    Тогда справедливо следующее:
    \begin{enumerate}
        \item $ W(t) = \det X(t)$
        \item $\dot{X} = A(t)X$
    \end{enumerate}
    Обозначим через $\xi_1(t), \ldots \xi_n(t) $ --- строки матрицы $X(t)$, то есть $ X(t) = \begin{pmatrix}
        \xi_1(t)\\
        \vdots\\
        \xi_n(t)
    \end{pmatrix}$
    $\Rightarrow$ уравнение $\dot{X} = AX$ можно переписать в виде:
    \[ \dot{\xi}_k = a^k X = \sum_{s=1}^{n} a_{ks} \xi_s \text{ , где } a^k= (a_{k1}, \ldots, a_{kn}) \text{ --- k-ая строка матрицы } A(t), k=1,\ldots, n\] 
    Используя факт \ref{fact:2.2}, получим:
    \[ \dot{W} = (\det X)^{\boldsymbol{\cdot}} =\begin{vmatrix}
        \dot{\xi_1(t)} \\
        \vdots \\
        {\xi_n(t)}
    \end{vmatrix}
    + \ldots + \begin{vmatrix}
        \xi_1(t)\\
        \vdots\\
        \xi_{n-1}(t)\\
        \dot{\xi_n(t)}
    \end{vmatrix} =
    \begin{vmatrix} \displaystyle
        \sum_{s=1}^{n} a_{1s}\xi_s\\
        \xi_2\\
        \vdots\\
        \xi_n
    \end{vmatrix} +\ldots + 
    \begin{vmatrix} 
        \xi_1\\
        \vdots\\
        \xi_{n-1}\\
        \displaystyle \sum_{s=1}^{n} a_{ns}\xi_s
    \end{vmatrix} \]    
    Воспользуемся фактом \ref{fact:2.3} и получим:
    \[ = a_{11}\cdot W + a_{22}\cdot W + \ldots + a_{nn}\cdot W = W\cdot trA  \]
    И, наконец, это был end of the proof. 
\end{proof}

\begin{fact}\label{fact:2.2}
    Для доказательства предыдущей леммы мы использовали этот факт:
    \[ \left( \det \begin{pmatrix}
        b_1(t) \\
        \vdots \\
        b_n(t)
    \end{pmatrix} \right)^{\boldsymbol{\cdot}}
    =  \begin{vmatrix}
        \dot{b_1(t)} \\
        \vdots \\
        {b_n(t)}
    \end{vmatrix}
    + \ldots + \begin{vmatrix}
        b_1(t)\\
        \vdots\\
        b_{n-1}(t)\\
        \dot{b_n(t)}
    \end{vmatrix}  \]
\end{fact}


\begin{fact}\label{fact:2.3}
    \[ \begin{vmatrix}
        a_1^{(1)} + a_1^{(2)} \\
        a_2 \\
        \vdots \\
        a_n
    \end{vmatrix} = \begin{vmatrix}
        a_1^{(1)} \\
        a_2 \\
        \vdots \\
        a_n
    \end{vmatrix} + \begin{vmatrix}
        a_1^{(2)} \\
        a_2 \\
        \vdots \\
        a_n \end{vmatrix}
        \]
\end{fact}

\section{Вронскиан скалярных функций}
Пусть $\mathcal{F} = \{ f: I\to \mathbb{R} | f\in C^{n-1}(I) \}, n\ge 1, n\in \mathbb{N}$.
\begin{definition}
    Пусть $f_1, \ldots f_n\in \mathcal{F}$. 
    \[ W_{f_1,\ldots, f_n}(t) = \det \begin{pmatrix}
        f_1(t) & f_2(t) & \ldots & f_n(t)\\
        f_1'(t) & f_2'(t) & \ldots & f_n'(t)\\
        \vdots & \vdots & \vdots & \vdots \\
        f_1^{(n-1)}(t)& f_2^{(n-1)}(t) & \ldots & f_n^{(n-1)}(t) 
    \end{pmatrix} \text{ --- вронскиан скалярных функций } f_1, \ldots, f_n \]
\end{definition}

\begin{statement}
    Если $f_1, \ldots, f_n$ --- линейно зависимые $\Rightarrow W_{f_1, \ldots, f_n}(t)  \equiv 0 \ \forall t\in I$.\\
    В обратную сторону это утверждение неверно.
\end{statement}

Рассматриваем ДУ 
\begin{equation}\label{eq:2}
    y^{(n)} + a_1(t)y^{(n-1)} + \ldots + a_{n-1}(t)y' + a_n(t)y=0, \ a_1(t), \ldots, a_n(t) \in C(I)
\end{equation}

\begin{theorem}
    Пусть $y_1(t), \ldots, y_n(t)$ --- решения уравнения \eqref{eq:2}. Тогда следующие утверждения эквиваленты:
    \begin{enumerate}
        \item $y_1(t), \ldots, y_n(t)$ --- линейно зависимы на $I$
        \item $W_{y_1, \ldots, y_n}(t) \equiv 0$ на $I$
        \item $\exists t_0\in I \Rightarrow W_{y_1, \ldots, y_n} (t_0) = 0$
    \end{enumerate}
\end{theorem}

\begin{remark}
    Уравнение \eqref{eq:2} можно свести к системе вида \eqref{eq:1}: $\dot{x} = A(t)x$, \\
    где $x(t) = (y(t), y'(t), \ldots, y^{(n-1)}(t))^{\text{T}}$, 
    \[ A(t) = \begin{pmatrix}
        0 & 1 & 0 & 0 & \dots & 0 \\
        0 & 0 & 1 & 0 & \dots & 0 \\
        \vdots & \vdots & \vdots & \vdots & \ddots & \vdots \\
        0 & 0 & 0 & 0 & \dots & 1 \\
        a_n & a_{n-1} & \dots & \dots & a_2 & a_1 \\
    \end{pmatrix} \]
    Диагональю в этой матрице является:
    \[
    A(t) =
    \begin{pmatrix}
        \textcolor{red}{0} & 1 & 0 & \cdots & 0 \\
        0 & \textcolor{red}{0} & 1 & \cdots & 0 \\
        \vdots & \vdots & \textcolor{red}{\ddots} & \ddots & \vdots \\
        0 & 0 & \cdots & \textcolor{red}{0} & 1 \\
        -a_n & -a_{n-1} & \cdots & -a_2 & \textcolor{red}{-a_1}
    \end{pmatrix} \]
    Также, если $y_1, \ldots, y_n$ --- решения уравнения \eqref{eq:2}, то 
    \[ W_{y_1, \ldots, y_n} (t) = W_{x_1, \ldots x_n}(t) \text{ , где } x_k(t) = \begin{pmatrix}
        y_k(t)\\
        y_k'(t)\\
        \vdots\\
        y_k^{(n-1)}(t) 
        
    \end{pmatrix}\] 
\end{remark}

\begin{theorem}
    (Лиувилля-Остроградского) Пусть $y_1(t), \ldots, y_n(t)$ --- решения уравнения \eqref{eq:2}, а $W$ --- вронскиан этого набора функций. Тогда
    \[  W(t) = W(t_0) e^{-\int_{t_0}^{t} a_1(s)\, ds} \ \forall t,t_0 \in I \] 
\end{theorem}


\end{document} 